\subsection{CONJUGACY OF CARTAN SUBALGEBRAS}

Let g be a Lie algebra, h a nilpotent subalgebra of g and R the set of non-zero weights of h in g, in other words the set of linear forms \(\lambda \neq 0\) on h such that \(\mathfrak{g}^{\lambda}(\mathfrak{h}) \neq 0\) , cf.~§1, no. 3, Prop. 9 (iii). Assume that

\[
\mathfrak {g} = \mathfrak {g} ^ {0} (\mathfrak {h}) \oplus \sum_ {\lambda \in R} \mathfrak {g} ^ {\lambda} (\mathfrak {h}),
\]

which is the case if \(k\) is algebraically closed (§1, no. 3, Prop. 9 (i)). For \(\lambda \in \mathbb{R}\) and \(x \in \mathfrak{g}^{\lambda}(\mathfrak{h})\) , ad \(x\) is nilpotent (§1, no. 3, Prop. 10 (iv)). Denote by \(\mathrm{E}(\mathfrak{h})\) the subgroup of \(\mathrm{Aut}_e(\mathfrak{g})\) generated by the \(e^{\mathrm{ad}x}\) where \(x\) is of the form above. If \(u \in \mathrm{Aut}(\mathfrak{g})\) , it is immediate that \(u\mathrm{E}(\mathfrak{h})u^{-1} = \mathrm{E}(u(\mathfrak{h}))\) .

Lemma 2. (i) Let \(\mathfrak{h}_r\) be the set of \(x\in \mathfrak{h}\) such that \(\mathfrak{g}^0 (x) = \mathfrak{g}^0 (\mathfrak{h})\) ; this is the set of \(x\in \mathfrak{h}\) such that \(\lambda (x)\neq 0\) for all \(\lambda \in \mathrm{R}\) , and \(\mathfrak{h}_r\) is open and dense in \(\mathfrak{h}\) in the Zariski topology.

\begin{enumerate}
\def\labelenumi{(\roman{enumi})}
\setcounter{enumi}{1}
\tightlist
\item
  Put \(R = \{\lambda_{1}, \lambda_{2}, \ldots, \lambda_{p}\}\) where the \(\lambda_{i}\) are mutually distinct. Let F be the map from \(\mathfrak{g}^{0}(\mathfrak{h}) \times \mathfrak{g}^{\lambda_{1}}(\mathfrak{h}) \times \cdots \times \mathfrak{g}^{\lambda_{p}}(\mathfrak{h})\) to g defined by the formula
\end{enumerate}

\[
\mathrm{F} (h, x _ {1}, \dots , x _ {p}) = e ^ {\operatorname{ad} x _ {1}} \dots e ^ {\operatorname{ad} x _ {p}} h.
\]

Then \(\mathrm{F}\) is a dominant polynomial map (App. I).

Assertion (i) is clear. We prove (ii). Let \(n = \dim \mathfrak{g}\) . If \(\lambda \in \mathrm{R}\) and \(x \in \mathfrak{g}^{\lambda}(\mathfrak{h})\) , we have \((\operatorname{ad} x)^n = 0\) . It follows that \((y, x) \mapsto e^{\operatorname{ad} x} y\) is a polynomial map from \(\mathfrak{g} \times \mathfrak{g}^{\lambda}(\mathfrak{h})\) to \(\mathfrak{g}\) ; it follows by induction that \(\mathrm{F}\) is polynomial. Let \(h_0 \in \mathfrak{h}_r\) and let \(\mathrm{DF}\) be the tangent linear map of \(\mathrm{F}\) at \((h_0, 0, \ldots, 0)\) ; we show that \(\mathrm{DF}\) is surjective. For \(h \in \mathfrak{g}^{0}(\mathfrak{h})\) , \(F(h_{0} + h, 0, \ldots, 0) = h_{0} + h\) , so \(\mathrm{DF}(h, 0, \ldots, 0) = h\) and \(\mathrm{Im}(\mathrm{DF}) \supset \mathfrak{g}^{0}(\mathfrak{h})\) . On the other hand, for \(x \in \mathfrak{g}^{\lambda_{1}}(\mathfrak{h})\) ,

\[
\mathrm{F} (h _ {0}, x, 0, \dots , 0) = e ^ {\operatorname{ad} x} h _ {0} = h _ {0} + (\operatorname{ad} x). h _ {0} + \frac {(\operatorname{ad} x) ^ {2}}{2 !} h _ {0} + \dots
\]

so \(\mathrm{DF}(0,x,0,\ldots ,0) = (\mathrm{ad}x).h_0 = -(\mathrm{ad}h_0)x\) ; since \(\mathrm{ad}h_0\) induces an automorphism of \(\mathfrak{g}^{\lambda_1}(\mathfrak{h})\) , \(\operatorname {Im}(\mathrm{DF})\supset \mathfrak{g}^{\lambda_1}(\mathfrak{h})\) . Similarly,

\[
\operatorname{Im} (\mathrm{DF}) \supset \mathfrak {g} ^ {\lambda_ {i}} (\mathfrak {h})
\]

for all \(i\) , hence the surjectivity of DF. Prop. 4 of App. I now shows that F is dominant.

PROPOSITION 2. Assume that \(k\) is algebraically closed. Let \(\mathfrak{g}\) be a Lie algebra, \(\mathfrak{h}\) and \(\mathfrak{h}'\) Cartan subalgebras of \(\mathfrak{g}\) . There exist \(u \in \mathrm{E}(\mathfrak{h})\) and \(u' \in \mathrm{E}(\mathfrak{h}')\) such that \(u(\mathfrak{h}) = u'(\mathfrak{h}')\) .

We retain the notation of Lemma 2. From the fact that \(\mathfrak{h}\) and \(\mathfrak{h}'\) are Cartan subalgebras, it follows that \(\mathfrak{g}^0 (\mathfrak{h}) = \mathfrak{h}\) and \(\mathfrak{g}^0 (\mathfrak{h}') = \mathfrak{h}'\) . By Lemma 2 and Prop. 3 of App. I, \(\mathrm{E}(\mathfrak{h})\mathfrak{h}_r\) and \(\mathrm{E}(\mathfrak{h}')\mathfrak{h}_r'\) contain open dense subsets of \(\mathfrak{g}\) in the Zariski topology. Thus \(\mathrm{E}(\mathfrak{h})\mathfrak{h}_r\cap \mathrm{E}(\mathfrak{h}')\mathfrak{h}_r' \neq \varnothing\) . In other words, there exist \(u\in \mathrm{E}(\mathfrak{h}), u' \in \mathrm{E}(\mathfrak{h}')\) , \(h\in \mathfrak{h}_r\) , \(h' \in \mathfrak{h}_r'\) such that \(u(h) = u'(h')\) ; then

\[
u (\mathfrak {h}) = u \left(\mathfrak {g} ^ {0} (\mathfrak {h})\right) = \mathfrak {g} ^ {0} (u (h)) = \mathfrak {g} ^ {0} \left(u ^ {\prime} \left(h ^ {\prime}\right)\right) = u ^ {\prime} \left(\mathfrak {h} ^ {\prime}\right).
\]

COROLLARY. \(\mathrm{E}(\mathfrak{h}) = \mathrm{E}(\mathfrak{h}')\) .

Let \(u, u'\) be as in Prop. 2. Then

\[
\mathrm{E} (\mathfrak {h}) = u \mathrm{E} (\mathfrak {h}) u ^ {- 1} = \mathrm{E} (u (\mathfrak {h})) = \mathrm{E} (u ^ {\prime} (\mathfrak {h} ^ {\prime})) = u ^ {\prime} \mathrm{E} (\mathfrak {h} ^ {\prime}) u ^ {- 1} = \mathrm{E} (\mathfrak {h} ^ {\prime}),
\]

hence the corollary.

Because of this result, if k is algebraically closed we shall denote simply by E the group \(\mathrm{E}(\mathfrak{h})\) , where h is a Cartan subalgebra of g.

In general, \(\operatorname{Aut}_{e}(\mathfrak{g}) \neq \operatorname{E}\) (for example, if g is nilpotent, E reduces to the identity element, even though non-trivial elementary automorphisms exist provided g is non-commutative). However, it can be shown (Chap. VIII, §10, Exerc. 5) that \(\operatorname{Aut}_{e}(\mathfrak{g}) = \operatorname{E}\) for g semi-simple.

THEOREM 1. Assume that k is algebraically closed. Let g be a Lie algebra. The group E is normal in \(\operatorname{Aut}(\mathfrak{g})\) and operates transitively on the set of Cartan subalgebras of g.

Let \(\mathfrak{h}\) be a Cartan subalgebra of \(\mathfrak{g}\) , and \(v\in\mathrm{Aut}(\mathfrak{g})\) . Then

\[
v \mathrm{E} (\mathfrak {h}) v ^ {- 1} = \mathrm{E} (v (\mathfrak {h})) = \mathrm{E} (\mathfrak {h}),
\]

so \(\mathrm{E}(\mathfrak{h}) = \mathrm{E}\) is normal in \(\mathrm{Aut}(\mathfrak{g})\) . If \(h'\) is another Cartan subalgebra of g, then, in the notation of Prop. 2, \(u'^{-1}u(\mathfrak{h}) = \mathfrak{h}'\) , and \(u'^{-1}u \in E\) .

